% 第十四章：其他方法简介

\chapter{其他方法简介}

\begin{wulibj}[本章导读]
除了前面介绍的主要方法外，数学物理方程还有许多其他求解方法。本章简要介绍变分法、摄动法和数值方法的基本思想，为后续学习和研究打下基础。
\end{wulibj}

% ========== 第一节 ==========

\subsection{变分原理}

许多物理问题可以表述为某个泛函的极值问题。例如：
\begin{itemize}
    \item 力学中的最小作用量原理
    \item 光学中的费马原理
    \item 静电学中的汤姆逊原理
\end{itemize}

\begin{dingliyi}{}{泛函}
设 $J[y]$ 是函数 $y(x)$ 的“函数”，称为\textbf{泛函}。变分法研究泛函的极值问题。
\end{dingliyi}

\subsection{欧拉-拉格朗日方程}

考虑泛函
\begin{equation}
    J[y] = \int_a^b F(x, y, y')\,\mathrm{d}x
\end{equation}

\begin{zhongyao}[欧拉-拉格朗日方程]
使泛函 $J[y]$ 取极值的函数 $y(x)$ 满足
\[
    \frac{\partial F}{\partial y} - \frac{\mathrm{d}}{\mathrm{d}x}\frac{\partial F}{\partial y'} = 0
\]
\end{zhongyao}

\begin{proof}
设 $y(x)$ 是极值函数，考虑变分 $y + \varepsilon\eta(x)$，其中 $\eta(a) = \eta(b) = 0$。

令
\begin{equation}
    \Phi(\varepsilon) = J[y + \varepsilon\eta] = \int_a^b F(x, y + \varepsilon\eta, y' + \varepsilon\eta')\,\mathrm{d}x
\end{equation}

极值条件：$\Phi'(0) = 0$。

\begin{equation}
    \Phi'(0) = \int_a^b \left(\frac{\partial F}{\partial y}\eta + \frac{\partial F}{\partial y'}\eta'\right)\,\mathrm{d}x = 0
\end{equation}

分部积分第二项：
\begin{equation}
    \int_a^b \frac{\partial F}{\partial y'}\eta'\,\mathrm{d}x = \left[\frac{\partial F}{\partial y'}\eta\right]_a^b - \int_a^b \frac{\mathrm{d}}{\mathrm{d}x}\frac{\partial F}{\partial y'}\eta\,\mathrm{d}x = -\int_a^b \frac{\mathrm{d}}{\mathrm{d}x}\frac{\partial F}{\partial y'}\eta\,\mathrm{d}x
\end{equation}

所以
\begin{equation}
    \int_a^b \left(\frac{\partial F}{\partial y} - \frac{\mathrm{d}}{\mathrm{d}x}\frac{\partial F}{\partial y'}\right)\eta\,\mathrm{d}x = 0
\end{equation}

由于 $\eta$ 任意，得到欧拉-拉格朗日方程。
\end{proof}

\subsection{变分法与微分方程的关系}

许多微分方程可以转化为变分问题。

\begin{lizi}{}{}
证明拉普拉斯方程 $\nabla^2 u = 0$ 对应于泛函 $J[u] = \frac{1}{2}\int_\Omega |\nabla u|^2\,\mathrm{d}V$ 的极小值问题。

\begin{jie}
    $F = \frac{1}{2}|\nabla u|^2 = \frac{1}{2}(u_x^2 + u_y^2 + u_z^2)$。
    
    欧拉-拉格朗日方程：
    \begin{equation}
        \frac{\partial F}{\partial u} - \sum_{i}\frac{\partial}{\partial x_i}\frac{\partial F}{\partial u_{x_i}} = 0 - \nabla^2 u = 0
    \end{equation}
    
    即 $\nabla^2 u = 0$。
    
    物理意义：调和函数使``能量泛函''最小。
\end{jie}
\end{lizi}

\subsection{里茨方法}

变分法可以用于数值求解：选择一组基函数 $\{\phi_n\}$，设近似解 $u \approx \sum_{n=1}^N c_n\phi_n$，代入泛函 $J[u]$，对系数 $c_n$ 求极值，得到线性方程组。

% ========== 第二节 ==========
\section{摄动法}

\subsection{正则摄动}

当方程含有小参数 $\varepsilon$ 时，可以将解展开为 $\varepsilon$ 的幂级数。

\begin{lizi}{}{}
求解摄动方程：
\begin{equation}
    y'' + y = \varepsilon y^2, \quad y(0) = 1, \quad y'(0) = 0
\end{equation}

\begin{jie}
    设 $y = y_0 + \varepsilon y_1 + \varepsilon^2 y_2 + \cdots$。
    
    代入方程，比较 $\varepsilon$ 的各阶系数：
    
    $O(1)$：
    \begin{equation}
        y_0'' + y_0 = 0, \quad y_0(0) = 1, \quad y_0'(0) = 0
    \end{equation}
    解得 $y_0 = \cos x$。
    
    $O(\varepsilon)$：
    \begin{equation}
        y_1'' + y_1 = y_0^2 = \cos^2 x, \quad y_1(0) = y_1'(0) = 0
    \end{equation}
    
    利用 $\cos^2 x = \frac{1 + \cos 2x}{2}$：
    \begin{equation}
        y_1'' + y_1 = \frac{1}{2} + \frac{1}{2}\cos 2x
    \end{equation}
    
    特解：$y_{1p} = \frac{1}{2} - \frac{1}{6}\cos 2x$。
    
    通解：$y_1 = A\cos x + B\sin x + \frac{1}{2} - \frac{1}{6}\cos 2x$。
    
    由初始条件得 $A = -\frac{1}{3}$，$B = 0$。
    
    所以
    \begin{equation}
        y \approx \cos x + \varepsilon\left(\frac{1}{2} - \frac{1}{3}\cos x - \frac{1}{6}\cos 2x\right)
    \end{equation}
\end{jie}
\end{lizi}

\subsection{奇异摄动}

当小参数出现在最高阶导数前时，正则摄动可能失效。这时需要使用\textbf{边界层方法}。

\begin{jinshi}[边界层现象]
当最高阶导数系数为小参数时，解可能在边界附近发生剧烈变化，形成``边界层''。正则展开无法捕捉这种行为。
\end{jinshi}

\subsection{多尺度方法}

对于长期项（随时间增长的项），可以使用多尺度方法：引入多个时间尺度 $T_0 = t$，$T_1 = \varepsilon t$，$T_2 = \varepsilon^2 t$，$\ldots$。

% ========== 第三节 ==========
\section{数值方法简介}

\subsection{有限差分法}

将偏导数用差商近似：
\begin{align}
    \frac{\partial u}{\partial x} &\approx \frac{u_{i+1,j} - u_{i-1,j}}{2\Delta x} \\
    \frac{\partial^2 u}{\partial x^2} &\approx \frac{u_{i+1,j} - 2u_{i,j} + u_{i-1,j}}{(\Delta x)^2}
\end{align}

将微分方程离散化为代数方程组。

\begin{figure}[htbp]
\centering
\begin{tikzpicture}[scale=0.8]
    % 网格
    \draw[gray, thin] (0, 0) grid (5, 4);
    
    % 点
    \foreach \i in {0,...,5} {
        \foreach \j in {0,...,4} {
            \fill (\i, \j) circle (1.5pt);
        }
    }
    
    % 中心点及其邻居
    \fill[red] (2.5, 2) circle (3pt);
    \fill[blue] (1.5, 2) circle (3pt);
    \fill[blue] (3.5, 2) circle (3pt);
    \fill[blue] (2.5, 1) circle (3pt);
    \fill[blue] (2.5, 3) circle (3pt);
    
    \node[red, below] at (2.5, 1.8) {$u_{i,j}$};
    \node[blue, left] at (1.3, 2) {$u_{i-1,j}$};
    \node[blue, right] at (3.7, 2) {$u_{i+1,j}$};
    \node[blue, below] at (2.5, 0.8) {$u_{i,j-1}$};
    \node[blue, above] at (2.5, 3.2) {$u_{i,j+1}$};
    
    \node at (2.5, -0.7) {五点差分格式};
\end{tikzpicture}
\caption{有限差分网格}
\label{fig:finite-difference}
\end{figure}

\subsection{有限元法}

将求解区域划分为有限个小单元，在每个单元内用插值函数近似解，通过变分原理得到离散方程。

基本步骤：
\begin{enumerate}
    \item 区域剖分
    \item 选择插值函数
    \item 单元分析
    \item 整体合成
    \item 施加边界条件
    \item 求解线性方程组
\end{enumerate}

\subsection{有限体积法}

基于积分形式的守恒定律，对每个控制体积建立离散方程。特别适合流体力学问题。

% ========== 第四节 ==========
\section{方法总结与选择}

\begin{table}[htbp]
\centering
\begin{tabular}{llll}
\toprule
\textbf{方法} & \textbf{适用方程类型} & \textbf{优点} & \textbf{缺点} \\
\midrule
行波法 & 波动方程 & 解析解，物理意义明确 & 只适用于初值问题 \\
分离变量法\index{分离变量法} & 有界区域问题 & 理论完整 & 只适用于规则区域 \\
积分变换法 & 无界/半无界问题 & 自动处理无穷远条件 & 逆变换困难 \\
格林函数\index{格林函数}法 & 线性问题 & 解的结构清晰 & 格林函数\index{格林函数}难求 \\
变分法 & 椭圆型问题 & 理论基础坚实 & 需要找到正确泛函 \\
摄动法 & 含小参数问题 & 渐近展开 & 可能出现长期项 \\
数值方法 & 一般问题 & 适用范围广 & 精度依赖网格 \\
\bottomrule
\end{tabular}
\caption{各种方法的比较}
\end{table}

% ========== 本章小结 ==========
\section{本章小结}

\subsection{变分法}

\begin{itemize}
    \item 泛函极值 $\Rightarrow$ 欧拉-拉格朗日方程
    \item 微分方程 $\Leftrightarrow$ 变分问题
    \item 应用：里茨方法、伽辽金方法
\end{itemize}

\subsection{摄动法}

\begin{itemize}
    \item 正则摄动：幂级数展开
    \item 奇异摄动：边界层方法
    \item 多尺度方法：消除长期项
\end{itemize}

\subsection{数值方法}

\begin{itemize}
    \item 有限差分法：简单直观
    \item 有限元法：适应复杂边界
    \item 有限体积法：保持守恒性
\end{itemize}

% ========== 习题 ==========
\section{习题}

\textbf{基础题}

\begin{enumerate}
    \item 推导泛函 $J[y] = \int_a^b (y'^2 + y^2)\,\mathrm{d}x$ 的欧拉-拉格朗日方程。
    
    \item 用正则摄动法求解 $y' + y = \varepsilon y^2$，$y(0) = 1$。
\end{enumerate}

\textbf{综合题}

\begin{enumerate}[resume]
    \item 证明：对于泛函 $J[u] = \int_\Omega \left[\frac{1}{2}|\nabla u|^2 - fu\right]\,\mathrm{d}V$，其欧拉-拉格朗日方程是泊松方程 $-\nabla^2 u = f$。
    
    \item 用有限差分法写出泊松方程 $u_{xx} + u_{yy} = f$ 在正方形网格上的五点差分格式。
\end{enumerate}

\textbf{挑战题}

\begin{enumerate}[resume]
    \item 研究奇异摄动问题 $\varepsilon y'' + y' + y = 0$，$y(0) = 0$，$y(1) = 1$ 的边界层结构。
\end{enumerate}
